\documentclass[12pt]{article}
\usepackage[utf8]{inputenc}
\usepackage[T1]{fontenc}
\usepackage{hyperref}
\title{Haruai Language Resources for K-12 Curriculum Development}
\author{Compiled by Claude Sonnet 5 (claude-sonnet-5) for Andr\'as Kornai}
\date{}
\begin{document}
\maketitle

\section{Introduction}
Haruai (ISO 639-3 tmd, Glottocode haru1245) is a Piawi-family language of Madang Province, Papua New Guinea. It is also called Waibuk, Wovan or Taman. Search results give about 2,000 speakers and three dialects (North/Hamil, Central/Mambar, South/Arama). One catalogue lists a much lower figure and marks the language endangered. The search found no school materials, primers or native-language texts. The material below is background and linguistic reference only.

\section{Primary Educational Resources}
No primers, textbooks or school materials in Haruai were found in the searches.

\section{Linguistic and Reference Materials}
\begin{itemize}
\item Haruai language (Wikipedia overview): \url{https://en.wikipedia.org/wiki/Haruai_language}
\item Piawi languages (Wikipedia): \url{https://en.wikipedia.org/wiki/Piawi_languages}
\item A Grammar of Haruai, project page at the Max Planck Institute for Evolutionary Anthropology. It describes a grammar by Bernard Comrie and John Davies (SIL). I could not confirm that it has been published. \url{https://www.eva.mpg.de/linguistics/past-research-resources/documentation-and-description/a-grammar-of-haruai/}
\item Grambank entry for Haruai, which includes Comrie's bibliography: \url{https://grambank.clld.org/languages/haru1245}
\item Haruai numerals data: \url{https://lingweb.eva.mpg.de/channumerals/Haruai.htm}
\item Glottolog entry for tmd: \url{https://glottolog.org/resource/languoid/iso/tmd}
\item WALS entry: \url{https://wals.info/languoid/lect/wals_code_har}
\end{itemize}
No Haruai dictionary appeared in the results.

\section{Religious and Community Resources}
None were found in the searches.

\section{Digital and Audio Resources}
No audio, video or software resources were found. The Verbix page below is a general endangered-language listing: \url{https://docs.verbix.com/EndangeredLanguages/Haruai}

\section{Curriculum Development Context}
Little usable material exists in the search results. Any curriculum would first need a documentation phase: a dictionary, a practical orthography and recorded texts. These could build on the Comrie and Davies grammar work, once its status and availability are confirmed with the authors or SIL.
\end{document}